\documentclass[11pt]{article}

\usepackage{amsmath,amsthm,amssymb}
\usepackage{mathtools}
\usepackage{booktabs}
\usepackage[margin=1in]{geometry}

\newtheorem{theorem}{Theorem}
\newtheorem{lemma}[theorem]{Lemma}
\newtheorem{proposition}[theorem]{Proposition}
\newtheorem{corollary}[theorem]{Corollary}
\newtheorem{definition}[theorem]{Definition}
\newtheorem{counterexample}[theorem]{Counterexample}

\theoremstyle{remark}
\newtheorem{remark}[theorem]{Remark}

\newcommand{\CC}{\mathbb{C}}
\newcommand{\ZZ}{\mathbb{Z}}
\newcommand{\ket}[1]{|#1\rangle}

\title{The $\mathfrak{sl}_n$ Cactus Representation Theorem:\\
  Interval Reversals, Crystal Structures, and a Corrected Statement}
\author{Clio}
\date{April 10, 2026}

\begin{document}
\maketitle

\begin{abstract}
We prove that the interval reversal operators on $V^{\otimes k}$
($V = \CC^n$, arbitrary $n \geq 2$), constructed from the swap
operator $P$ (the ice-rule $R$-matrix at $q = 0$), define a
representation of the cactus group~$J_k$. The proof is dimension-free:
it depends only on the fact that $P$ is the permutation matrix on
$V \otimes V$, not on the dimension of~$V$. This extends the
$\mathfrak{sl}_2$ result to all $\mathfrak{sl}_n$.

We also analyse the relationship between this representation and the
Henriques--Kamnitzer crystal commutor, finding that they are
\emph{distinct}: the reversal permutation does not commute with
Kashiwara crystal operators on a single tensor product. The crystal
commutor for identical factors $B(\omega_1)^{\otimes k}$ is trivial
(the identity), while the reversal is nontrivial. The reversal does,
however, intertwine the standard and reordered tensor product crystal
structures---a tautological consequence of the signature rule being
defined relative to a reading order.
\end{abstract}

%% -----------------------------------------------------------------------
\section{The Theorem}
%% -----------------------------------------------------------------------

\begin{theorem}[$\mathfrak{sl}_n$ Cactus Representation]\label{thm:main}
Let $n \geq 2$, $V = \CC^n$ with standard basis
$\{e_1, \ldots, e_n\}$, and $P \in \mathrm{End}(V \otimes V)$ the
swap operator: $P(e_a \otimes e_b) = e_b \otimes e_a$. This is the
six-vertex ice-rule $R$-matrix at $q = 0$ for~$\mathfrak{sl}_n$.

For $k \geq 2$ and $1 \leq i < j \leq k$, define the commutor
$\sigma_{[i,j]} \in \mathrm{End}(V^{\otimes k})$ as the operator that
reverses the order of tensor factors $i$ through~$j$, acting as the
identity on all other factors. Then:
\begin{enumerate}
  \item \textbf{Cactus representation.}
    The map $\rho\colon J_k \to \mathrm{GL}(V^{\otimes k})$ defined by
    $s_{[i,j]} \mapsto \sigma_{[i,j]}$ is a group homomorphism.
  \item \textbf{Factorisation through $S_k$.}
    The representation factors as $\rho = \rho_{S_k} \circ \phi$, where
    $\phi\colon J_k \twoheadrightarrow S_k$ is the natural surjection
    and $\rho_{S_k}$ is the permutation representation on tensor factors.
  \item \textbf{Dimension-free.}
    The proof uses only the symmetric group structure, not
    $\dim V$. It holds identically for all $n \geq 2$.
\end{enumerate}
\end{theorem}

%% -----------------------------------------------------------------------
\section{Proof of Theorem~\ref{thm:main}}
%% -----------------------------------------------------------------------

The proof is identical to the $\mathfrak{sl}_2$ case (see our earlier
proof of the Cactus Representation Theorem, April~10, 2026). We reproduce
it here for completeness, noting that no step uses $\dim V = 2$.

\subsection{Construction}

\begin{definition}
The commutor $\sigma_{[i,j]} \in \mathrm{End}(V^{\otimes k})$ is the
image of the longest element $w_0 \in S_{j-i+1}$ under the permutation
representation $S_{j-i+1} \to \mathrm{GL}(V^{\otimes(j-i+1)})$,
$s_r \mapsto P_{i+r-1}$, embedded into $\mathrm{End}(V^{\otimes k})$.
Explicitly:
\[
  \sigma_{[i,j]}\ket{b_1 \cdots b_k}
  = \ket{b_1 \cdots b_{i-1}\, b_j\, b_{j-1} \cdots b_{i+1}\, b_i\, b_{j+1} \cdots b_k},
\]
where $b_r \in \{1, \ldots, n\}$ for all~$r$.
\end{definition}

\subsection{Verification of cactus relations}

\subsubsection*{Relation (C1): Involutivity}

$\sigma_{[i,j]}$ is the image of $w_0 \in S_{j-i+1}$ under a group
homomorphism. Since $w_0^2 = e$ in the symmetric group, we have
$\sigma_{[i,j]}^2 = I$. Alternatively: reversing a substring twice
restores the original order. \qed

\subsubsection*{Relation (C2): Disjoint commutativity}

If $[i,j] \cap [k,l] = \emptyset$, then $\sigma_{[i,j]}$ and
$\sigma_{[k,l]}$ act on disjoint sets of tensor positions. Products of
operators $P_r$ with $r \in [i, j-1]$ commute with products of
$P_r$ with $r \in [k, l-1]$, since $P_r P_s = P_s P_r$ for
$|r - s| \geq 2$. \qed

\subsubsection*{Relation (C3): Nesting}

\begin{proposition}
If $[k,l] \subseteq [i,j]$, then
$\sigma_{[i,j]}\, \sigma_{[k,l]} = \sigma_{[i+j-l,\, i+j-k]}\, \sigma_{[i,j]}$.
\end{proposition}

\begin{proof}
Identify positions $\{i, \ldots, j\}$ with $\{1, \ldots, m\}$ via
$t \mapsto t - i + 1$, where $m = j - i + 1$. Then $\sigma_{[i,j]}$
corresponds to $w_0 \in S_m$ (the map $t \mapsto m + 1 - t$), and
$\sigma_{[k,l]}$ to $w_0^{[k',l']}$ (reversal on $[k', l']$ with
$k' = k - i + 1$, $l' = l - i + 1$).

The nesting relation becomes $w_0\, w_0^{[k',l']}\, w_0^{-1} = w_0^{[m+1-l',\, m+1-k']}$,
which we verify by computing the action on an arbitrary position $t$:

\textbf{Case 1:} $t \notin [m+1-l',\, m+1-k']$. Then
$w_0(t) = m+1-t \notin [k',l']$, so $w_0^{[k',l']}$ fixes it, giving
$\tau(t) = w_0^{-1}(w_0(t)) = t$.

\textbf{Case 2:} $t \in [m+1-l',\, m+1-k']$. Write $t = m+1-s$
with $s \in [k',l']$. Then $\tau(t) = w_0^{-1}(k'+l'-s)
= m+1-k'-l'+s = (m+1-k') + (m+1-l') - t$, which is the reversal
within $[m+1-l', m+1-k']$.
\end{proof}

\begin{remark}
This proof uses \emph{only} the symmetric group --- the fact that
$\sigma_{[i,j]}$ permutes tensor factor positions. It never references
the dimension of~$V$, the crystal structure, or any Lie-theoretic data.
The same proof works verbatim for $V = \CC^n$ (any $n \geq 1$), and
indeed for any vector space~$V$.
\end{remark}

%% -----------------------------------------------------------------------
\section{Relationship to the Henriques--Kamnitzer Crystal Commutor}
\label{sec:crystal}
%% -----------------------------------------------------------------------

The original conjecture (PROVE.md) asserted that the reversal
representation $\rho$ coincides with the Henriques--Kamnitzer crystal
commutor on $B(\omega_1)^{\otimes k}$ and intertwines the crystal
operators $\tilde{e}_\alpha$, $\tilde{f}_\alpha$. We show that both
claims are \textbf{false}, and identify the correct statement.

\subsection{The HK commutor on identical crystals is trivial}

\begin{proposition}\label{prop:hk-trivial}
Let $B = B(\omega_1)$ be the fundamental crystal of
$\mathfrak{sl}_n$. The Henriques--Kamnitzer crystal commutor
$\sigma_{\mathrm{HK}}$ on $B^{\otimes k}$ acts as the identity.
\end{proposition}

\begin{proof}
The HK commutor $\sigma_{B_1,B_2}\colon B_1 \otimes B_2 \to
B_2 \otimes B_1$ is the unique crystal isomorphism between the two
tensor product crystals. When $B_1 = B_2 = B$, the objects
$B_1 \otimes B_2$ and $B_2 \otimes B_1$ are \emph{equal} as crystals
--- the Kashiwara tensor product rule depends on the crystal data
$(\varphi_i, \varepsilon_i)$ of each element, not on which ``copy'' it
belongs to.

The tensor product $B(\omega_1) \otimes B(\omega_1)$ decomposes as
$B(2\omega_1) \oplus B(\omega_2)$ for $n \geq 3$ (and
$B(2\omega_1) \oplus B(0)$ for $n = 2$), each component with
multiplicity~1. The unique crystal automorphism of a connected crystal
$B(\lambda)$ is the identity (it is determined by the image of the
highest weight vector, which must be the unique highest weight vector).
Therefore $\sigma_{\mathrm{HK}}$ is the identity on each component,
hence the identity on $B \otimes B$.

By the coherence condition of coboundary categories, the commutor for
composite objects is built from pairwise commutors:
\[
  \sigma_{B^{\otimes p}, B^{\otimes q}} =
  (\sigma_{B, B^{\otimes q}} \otimes \mathrm{id}_{B^{\otimes(p-1)}})
  \circ \cdots
\]
When each pairwise commutor is the identity, all composite commutors
are the identity.
\end{proof}

\subsection{The reversal is not a crystal morphism}

\begin{counterexample}\label{ce:swap}
For $\mathfrak{sl}_2$, $B = B(\omega_1) = \{1, 2\}$ with
$\tilde{f}_1(1) = 2$. The tensor product $B \otimes B$ has crystal
structure (Kashiwara rule, left factor has priority when
$\varphi_1(\text{left}) > \varepsilon_1(\text{right})$):
\[
  \tilde{f}_1(1, 1) = (2, 1), \quad
  \tilde{f}_1(2, 1) = (2, 2), \quad
  \tilde{f}_1(1, 2) = 0.
\]
The swap $P\colon (a,b) \mapsto (b,a)$ fails to commute with
$\tilde{f}_1$:
\[
  P(\tilde{f}_1(1, 1)) = P(2, 1) = (1, 2), \qquad
  \tilde{f}_1(P(1, 1)) = \tilde{f}_1(1, 1) = (2, 1).
\]
Since $(1, 2) \neq (2, 1)$, the swap $P$ is \textbf{not} a crystal
endomorphism of $B \otimes B$. The same failure occurs for all
$n \geq 2$.
\end{counterexample}

\begin{corollary}
The reversal $\sigma_{[i,j]}$ does not commute with the Kashiwara
crystal operators $\tilde{f}_\alpha$ on
$(B(\omega_1)^{\otimes k}, \Delta_{\mathrm{std}})$.
\end{corollary}

\subsection{What the reversal does intertwine}

The reversal $\sigma_{[i,j]}$ is not a crystal endomorphism, but it
is a crystal \emph{isomorphism between different crystal structures}
on the same underlying set.

\begin{definition}
For a reading order $\pi = (\pi_1, \ldots, \pi_k)$ (a permutation of
$\{1, \ldots, k\}$), define the $\pi$-crystal structure on
$\{1, \ldots, n\}^k$ by the Kashiwara tensor product rule with
the $\alpha_i$-signature read in the order $\pi_1, \pi_2, \ldots, \pi_k$
instead of $1, 2, \ldots, k$.
\end{definition}

\begin{proposition}\label{prop:intertwine}
Let $\pi_{[p,q]}$ denote the reading order obtained from the standard
order $(1, \ldots, k)$ by reversing positions $p$ through~$q$. Then
the reversal $\sigma_{[p,q]}$ is a crystal isomorphism
\[
  \sigma_{[p,q]}\colon
  (B(\omega_1)^{\otimes k}, \Delta_{\mathrm{std}})
  \xrightarrow{\;\sim\;}
  (B(\omega_1)^{\otimes k}, \Delta_{\pi_{[p,q]}}).
\]
\end{proposition}

\begin{proof}
The Kashiwara tensor product rule determines $\tilde{f}_i(w)$ from the
$\alpha_i$-signature of the word $w = (w_1, \ldots, w_k)$, which is a
string of $+$ and $-$ symbols at positions where $w_j = i$ and
$w_j = i+1$ respectively, read in a specified order. The reduced
signature (after cancelling $+-$ pairs) determines which position $j_0$
is modified by $\tilde{f}_i$ (the leftmost uncancelled~$+$).

The reversal $\sigma_{[p,q]}$ simultaneously:
\begin{enumerate}
  \item permutes the values in positions $p, \ldots, q$, and
  \item corresponds to reading those positions in the reversed order.
\end{enumerate}
Since the signature symbols are determined by the \emph{values} at each
position, and the reduction algorithm depends on the \emph{order} in
which they are read, applying $\sigma_{[p,q]}$ to the values while
reversing the reading order produces exactly the same reduced signature
(with position labels relabeled). Therefore
\[
  \sigma_{[p,q]} \circ \tilde{f}_i^{\Delta_{\mathrm{std}}}
  = \tilde{f}_i^{\Delta_{\pi_{[p,q]}}} \circ \sigma_{[p,q]}.
\]
\end{proof}

\begin{remark}
This intertwining is essentially tautological: permuting both the values
\emph{and} the reading order preserves the reduced signature and hence
the crystal action. It does not give a crystal endomorphism (which would
require preserving a \emph{single} crystal structure).
\end{remark}

\subsection{The two representations compared}

\begin{center}
\renewcommand{\arraystretch}{1.3}
\begin{tabular}{@{}lcc@{}}
  \toprule
  Property & Reversal $\sigma_{[i,j]}$ & HK commutor $\sigma_{\mathrm{HK}}$ \\
  \midrule
  Action on $B^{\otimes k}$ & Nontrivial (reversal) & Trivial (identity) \\
  Crystal endomorphism? & No & Yes (trivially) \\
  Satisfies cactus relations? & Yes & Yes (trivially) \\
  Source & $R$-matrix at $q = 0$ & Coboundary category \\
  Intertwines & $\Delta_{\mathrm{std}} \leftrightarrow \Delta_{\pi}$ & $\Delta_{\mathrm{std}} \to \Delta_{\mathrm{std}}$ \\
  \bottomrule
\end{tabular}
\end{center}

\begin{remark}
The HK commutor becomes nontrivial for tensor products of
\emph{different} crystals: $\sigma_{B(\lambda), B(\mu)}$ with
$\lambda \neq \mu$ is the combinatorial $R$-matrix, which is a
nontrivial crystal isomorphism between $B(\lambda) \otimes B(\mu)$ and
$B(\mu) \otimes B(\lambda)$ (genuinely different crystal structures).
The triviality for identical crystals is a consequence of
$B \otimes B = B \otimes B$, not a deficiency of the theory.
\end{remark}

%% -----------------------------------------------------------------------
\section{Computational Verification}
%% -----------------------------------------------------------------------

\subsection{Cactus relations}

All three cactus relations (C1)--(C3) were verified by explicit
permutation matrix computation for the following cases:

\begin{center}
\begin{tabular}{@{}llrrrl@{}}
  \toprule
  $\mathfrak{g}$ & $\dim V$ & $k$ & $\dim V^{\otimes k}$ & Generators & Result \\
  \midrule
  $\mathfrak{sl}_2$ & 2 & 3, 4 & 8, 16 & 3, 6 & ALL PASS \\
  $\mathfrak{sl}_3$ & 3 & 3, 4 & 27, 81 & 3, 6 & ALL PASS \\
  $\mathfrak{sl}_4$ & 4 & 3, 4 & 64, 256 & 3, 6 & ALL PASS \\
  \bottomrule
\end{tabular}
\end{center}

\subsection{Crystal morphism failure}

The swap $P\colon (a,b) \mapsto (b,a)$ was verified to \emph{not}
commute with crystal operators $\tilde{f}_i$ for
$\mathfrak{sl}_n$, $n = 2, 3, 4$.

For $\mathfrak{sl}_2$, $k = 3$: all 5 non-zero crystal operator
applications fail to commute with the reversal.

For $\mathfrak{sl}_3$, $k = 3$: 22 out of 28 non-zero applications
fail (the 6 that ``pass'' do so by accident, not by structure).

\subsection{Crystal intertwining}

The intertwining property
$\sigma_{[p,q]} \circ \tilde{f}_i^{\Delta_{\mathrm{std}}}
= \tilde{f}_i^{\Delta_{\pi_{[p,q]}}} \circ \sigma_{[p,q]}$
was verified for all intervals $[p,q]$ with $\mathfrak{sl}_2$
($k = 3, 4$) and $\mathfrak{sl}_3$ ($k = 3, 4$). All tests pass.

%% -----------------------------------------------------------------------
\section{Discussion}
%% -----------------------------------------------------------------------

\subsection{What the theorem establishes}

The $\mathfrak{sl}_n$ Cactus Representation Theorem confirms that the
algebraic structure of interval reversals in the quasi-solvable regime
($q = 0$) is \emph{universal}: it is the cactus group $J_k$, regardless
of the rank $n$ of the Lie algebra. This extends the categorical phase
diagram from $\mathfrak{sl}_2$ to all $\mathfrak{sl}_n$:

\begin{center}
\renewcommand{\arraystretch}{1.3}
\begin{tabular}{@{}llll@{}}
  \toprule
  Regime & Category & Structure group & Holds for \\
  \midrule
  Solvable ($q \neq 0$) & Braided monoidal & $B_k$ & all $\mathfrak{sl}_n$ \\
  Quasi-solvable ($q = 0$) & Coboundary monoidal & $J_k$ & all $\mathfrak{sl}_n$ \\
  Symmetric ($q = 0, u = 0$) & Symmetric monoidal & $S_k$ & all $\mathfrak{sl}_n$ \\
  \bottomrule
\end{tabular}
\end{center}

\subsection{What the theorem does not establish}

The reversal representation is a \emph{group representation}, not a
\emph{crystal morphism}. It acts on the vector space $V^{\otimes k}$
by permuting basis vectors, but it does not preserve the crystal
structure defined by the Kashiwara operators. The crystal commutor of
Henriques--Kamnitzer acts trivially on $B(\omega_1)^{\otimes k}$ when
all factors are identical.

This suggests that the correct setting for the coboundary monoidal
structure at $q = 0$ is the category of \emph{representations}
(where the swap is the categorical braiding/commutor), not the category
of \emph{crystals} (where the commutor for identical objects is trivial).
This is consistent with the Alqady--Stroinski result that the
Temperley--Lieb category $\mathrm{TL}_0$ (a category of representations,
not crystals) is coboundary monoidal.

\subsection{Corrected relationship to Henriques--Kamnitzer}

The reversal permutation $\sigma_{[i,j]}$ and the HK crystal commutor
are related but distinct:

\begin{itemize}
  \item Both define cactus group representations on $B(\omega_1)^{\otimes k}$.
  \item The HK commutor is the identity (trivial representation).
  \item The reversal is the image of $J_k$ in $S_k$ acting by permuting
    tensor factor positions.
  \item The reversal intertwines different tensor product crystal
    structures (standard $\leftrightarrow$ reordered).
  \item The HK commutor is a crystal endomorphism (identity) of
    a single crystal structure.
\end{itemize}

The meaningful connection to HK arises for tensor products of
\emph{different} crystals $B(\lambda_1) \otimes \cdots \otimes B(\lambda_k)$,
where the HK commutor is nontrivial and the combinatorial $R$-matrix
provides the crystal isomorphism. The passage from the quantum group
$R$-matrix to the combinatorial $R$-matrix under $q \to 0$ is the
crystal-theoretic content of the braided $\to$ coboundary transition.

%% -----------------------------------------------------------------------
\section*{Acknowledgements}
%% -----------------------------------------------------------------------
Computational verification performed using NumPy and a custom
implementation of the Kashiwara tensor product rule (signature method).

%% -----------------------------------------------------------------------
\begin{thebibliography}{9}

\bibitem{henriques-kamnitzer}
A.~Henriques and J.~Kamnitzer,
\emph{Crystals and coboundary categories},
Duke Math.\ J.\ \textbf{132} (2006), no.~2, 191--216.

\bibitem{alqady-stroinski}
R.~Alqady and M.~Stroinski,
\emph{On the coboundary monoidal structure of the Temperley--Lieb category
at $q=0$},
preprint.

\bibitem{kashiwara}
M.~Kashiwara,
\emph{On crystal bases of the $q$-analogue of universal enveloping
algebras},
Duke Math.\ J.\ \textbf{63} (1991), no.~2, 465--516.

\bibitem{buciumas-scrimshaw}
V.~Buciumas and T.~Scrimshaw,
\emph{Quasi-solvable lattice models and Demazure atoms},
preprint.

\end{thebibliography}

\end{document}
